% Note for D. M. Gordon: a summary of run 1 (tag signed-difference-sets/run1).
% Build: latexmk (see latexmkrc); the PDF is build/main.pdf.
\documentclass[11pt]{amsart}

\usepackage{fontspec}
\setmonofont{DejaVu Sans Mono}[Scale=MatchLowercase]
\usepackage{amssymb}
\usepackage{booktabs}
\usepackage{tabularx}
\usepackage{longtable}
\usepackage{fancyvrb}
\usepackage[margin=2.5cm]{geometry}
\usepackage[hidelinks]{hyperref}

\input{build/numbers}
\emergencystretch=1.5em % room for lines with long code words
\newcommand{\zrows}[1]{\csname zrows#1\endcsname}
\newcommand{\repo}{\url{https://github.com/qedavid/solverz}}
\newcommand{\todo}[1]{\textbf{[TODO: #1]}}

\title[Open cases of the SDS table]{Settling open cases of the signed difference set table,
with Lean-checked proofs: a summary}
\author{David Schneider}
\address{Independent researcher, Austria\newline
ORCID: \href{https://orcid.org/0009-0009-4935-8285}{0009-0009-4935-8285}}
\email{david.schneider1920@gmail.com}
\date{\todo{date}}

\begin{document}

\begin{abstract}
Of the \openAfterHeWu{} cases of Gordon's table of signed difference sets~\cite{gordon-db} that
remain open after He and Wu~\cite{he-wu}, we settle \ourRows: \ourNone{} have no signed difference set and \ourExists{} have
one. A further \corrections{} cases listed as open are already settled by published
constructions; we list them as corrections.
Every case we settle or correct is checked by the Lean~4 proof assistant. The research, the proofs and
their Lean formalization were produced with AI tools under my direction; see the end of the
note. This note summarizes the results, the database corrections and
how to check them.
\end{abstract}

\maketitle

\section{Results}

A signed difference set (SDS) with parameters $(v,k,\lambda)$ in an abelian group $G$ of
order $v$ is a map $D \colon G \to \{0, \pm 1\}$ with $k$ nonzero values such that
$\sum_x D(x) D(x+g) = \lambda$ for every $g \neq 0$~\cite{gordon}. We write $n = k - \lambda$
and $s^2 = k + \lambda(v-1)$ for its squared sign sum, a square by~\cite[Lemma~1.1]{gordon}.
Most rules work with the \emph{coset sums} $a_h = \sum_{\phi(x) = h} D(x)$ of $D$ on a
quotient $\phi \colon G \to H$, with $c = v/|H|$ (the intersection numbers
of~\cite[Lemma~5.2]{gordon}); they satisfy $|a_h| \le c$ and $\sum_h a_h a_{h+g} = \lambda c$
for $g \neq 0$. Theorem numbers of cited papers are those of their arXiv versions, where there is one.

Table~\ref{tab:rules} lists the rules. Each is a Lean theorem that holds for every group (Z25
and Z29 are two each; Z1, Z2 and Z26 settle single cases directly); a row of the table is
settled by checking the rule's hypotheses for that row inside Lean. Rows are counted once, by
the first rule in the order the rules were found (their numbers). ``Ours'' means we found no
earlier statement for signed difference sets; it does not mean the idea is new.

Most rules share one shape. If an SDS existed, its coset sums on a suitable quotient $H$ would be
integers with $|a_h| \le c$, the sum and autocorrelation above, and often further conditions:
congruences from self-conjugate primes, values of characters, or invariance under multipliers.
Each rule shows that for its cases no integer vector meets them all. The exceptions work on $G$
itself: support parity (Z11) and the exhaustive searches of Z1, Z15 and Z27 (the latter two
over multiplier orbits); Z33 lifts the coset sums back to $G$.

\begingroup\small
\begin{longtable}{@{}l >{\raggedright\arraybackslash}p{0.5\textwidth} r >{\raggedright\arraybackslash}p{0.26\textwidth}@{}}
\caption{Rules settling open cases (no SDS).}\label{tab:rules}\\
\toprule
& Rule & Rows & Origin \\ \midrule
\endfirsthead
\toprule
& Rule & Rows & Origin \\ \midrule
\endhead
\multicolumn{4}{@{}l}{\emph{Congruences from self-conjugacy}} \\
Z8 & Self-conjugacy: $p^{2j} \mid n$, $p \nmid |H|$, $p$ self-conjugate modulo $\exp H$; the
  $a_h$ agree modulo $p^j$, impossible when $p^j > 2c$ & \zrows{8} & Turyn~\cite{turyn};
  SDS form ours (odd exponent:~\cite{he-wu}) \\
Z9 & Moments: the same congruence with $\sum a_h = \pm s$ and $\sum a_h^2 = k + \lambda(c-1)$
  & \zrows{9} & ours, on~\cite[Lemma~5.2]{gordon} \\
Z10 & Box bound: the same, with $|a_h| \le c$ & \zrows{10} & ours \\
Z22 & The congruence for several self-conjugate primes at once (CRT) & \zrows{22} & ours \\
\addlinespace
\multicolumn{4}{@{}l}{\emph{Counting and parity}} \\
Z11 & Support parity: $v$ odd, $\lambda$ even, $d(d-1) < v-1$ with $d = v - k$ & \zrows{11}
  & classical for $\lambda = 0$~\cite{wallis-whiteman}; general $\lambda$ ours (for $d = 1$,
  $v = p^m$, $p \equiv 3 \pmod 4$, $m$ even:~\cite[Cor.~4.5]{he-wu}) \\
Z16 & Marked elements per coset (an LP certificate) & \zrows{16} & ours \\
\addlinespace
\multicolumn{4}{@{}l}{\emph{Characters}} \\
Z17 & Real characters on a quotient $C_2^r$ & \zrows{17} & extends~\cite[Thm.~4.2]{he-wu} \\
Z18 & Characters on a quotient of exponent 4 & \zrows{18} & ours \\
Z19 & Characters on a quotient of exponent 3 & \zrows{19} & classical (the ramified prime above 3); SDS form ours \\
Z25 & Sylow quotient of order $p^t$: $p^{\min(t, \lceil e/2 \rceil)} \mid s$, $e = v_p(n)$;
  for $p = 2$ a sharper equality case & \zrows{25} & ours, on Turyn's argument \\
\addlinespace
\multicolumn{4}{@{}l}{\emph{Ma's lemma}} \\
Z24 & Ma's lemma on a quotient $C_{p^j} \times K$, $p$ self-conjugate on $K$:
  impossible when $p^{\lfloor e/2 \rfloor} > 2c$ & \zrows{24} & Ma's lemma~\cite{pott};
  SDS form ours \\
Z29 & Ma's lemma at its edge ($p = 2$, $4 \mid n$, on a quotient $C_{2^j} \times K$ as in Z24):
  $8 \mid 2s' - (k + \lambda(2c-1))$ with $s' = \sum D = \pm s$; for $c = 1$, a halved SDS &
  \zrows{29} & halving for $c = 1$: the argument of~\cite[Lemma~6.2]{schmidt-smith} (there for
  weight 36);
  general $c$ ours \\
Z32 & The halved coset sums for any $c$ (even contraction) & \zrows{32} & ours, on Ma's lemma \\
\addlinespace
\multicolumn{4}{@{}l}{\emph{Multipliers}} \\
Z15 & Multiplier theorem ($\gcd(v,n) = 1$, $n$ squarefree), then an exhaustive search over
  multiplier orbits & \zrows{15} & McFarland, as in~\cite[Thm.~5.1]{gordon};
  fixed translate~\cite{mcfarland-rice} \\
Z27 & The same for every $n$ prime to $v$ (Lean proof through Witt vectors) & \zrows{27} & as Z15 \\
Z28 & Multiplier-fixed coset sums on a quotient, searched exhaustively & \zrows{28} & the search of~\cite[\S5]{gordon}; its Lean proof ours \\
Z31 & Self-conjugate primes enlarge the multiplier group of the coset sums & \zrows{31} &
  for difference sets~\cite{gordon-schmidt,miller}; SDS form ours \\
Z30 & Field descent without ramification: $\varphi(m)\,n \le m F c^2$ on a quotient
  $\mathbb{Z}_m$, $\gcd(m,n) = 1$, $F \mid m$, and every $t \equiv 1 \pmod F$ is prime to $m$ and congruent
  modulo $m$ to a power of each prime dividing $n$ & \zrows{30} & Schmidt's field descent~\cite{tan};
  this form ours \\
\addlinespace
\multicolumn{4}{@{}l}{\emph{Exhaustive search}} \\
Z1, Z2, & Exhaustive search and contraction for & 4 & ours \\
Z26 & single small cases & & \\
Z33 & Contraction to a quotient, then a lift coset by coset & \zrows{33} & as in~\cite[\S5]{gordon}; the rows and the Lean proof ours \\
\bottomrule
\end{longtable}
\endgroup

The \ourExists{} cases with a set are listed in Appendix~\ref{app:sets}; Lean checks each set.
Its last column is the Solution's ID in the repository's \texttt{solutions.md}; each set is there
as data, in \path{lean/SignedDifferenceSets/Answer/N<case>.lean} and as a text file in
\path{objects/}.
They come from exhaustive and orbit searches, from cyclotomic classes, from partial difference
sets ($P - \{0\}$ is an SDS when $\lambda' - \mu = 2$, the step in the proof
of~\cite[Thm.~5.3]{he-wu}), and from signed combinations of relative difference sets (Bose's
affine and Pott's semiregular ones, see~\cite{pott}). One such combination gives, on paper, a
cyclic family $((q-1)(q^2+1),\, q^2+q-1,\, q-1)$ for every prime power $q$; Lean checks its two
members that were open, $q = 7$ and $q = 8$. The $(27,10,1)$ set in $C_3^3$ is the construction
of~\cite[Thm.~4.1]{he-chen-ge} at $m = 1$, outside its stated range.

For example, the $(25,12,1)$ set in $C_5 \times C_5$ (case N49), with $D(a,b)$ in row $a$
and column $b$, $a, b = 0, \dots, 4$, and blanks for $0$:
\[
\csname @@input\endcsname build/example % TeX's primitive \input, as for the appendix
\]

\section{Corrections to the database}

These \corrections{} cases are listed as open in the database~\cite{gordon-db},
although published constructions settle them. For each we give a set checked in Lean. We found
them while checking published constructions against the open cases; the list is not complete
(Section~\ref{sec:open}).

\begin{center}\small
\begin{tabular}{@{}lll@{}}
\toprule
Cases & Parameters & Construction \\ \midrule
N10864, N12056, N13963, & $(211,106,51)$, $(223,112,54)$, $(239,120,58)$, & \cite[Prop.~2.3]{gordon}
  (Paley and zero) \\
N16190, N18527, N19488, & $(251,126,61)$, $(263,132,64)$, $(271,136,66)$, & \\
N49190 & $(419,210,103)$, all cyclic & \\
N67 & $(27,14,5)$ in $C_3^3$ & \cite[Prop.~2.3]{gordon} in $\mathrm{GF}(27)$ \\
N73 & $(27,17,8)$ in $C_3^3$ & \cite[Thm.~4.1]{gordon}, $q = 3$ \\
N49624 & $(421,316,25)$ cyclic & \cite[Prop.~5.5(2)]{he-chen-ge} \\
N14471 & $(243,82,1)$ in $C_3^5$ & \cite[Thm.~4.1]{he-chen-ge} with $x_1 = 0$ \\
\bottomrule
\end{tabular}
\end{center}

For N14471, Theorem~4.1 of~\cite{he-chen-ge} needs $x_1 = 0$: for $x_1 \neq 0$, as in their
Example~4.2 ($x_1 = (1,0)$), the construction is not a signed set.

\section{How to check it}

The Lean library lives in the repository \repo, in the folder
\path{problems/signed-difference-sets/lean} (Lean~4 v4.35.0 with Mathlib). The statement every
result refers to is:

\begin{Verbatim}[fontsize=\small]
def IsSDS (k : ℕ) (lam : ℤ) (D : G → SignType) : Prop :=
  (Finset.univ.filter fun x => D x ≠ 0).card = k ∧
    ∀ g : G, g ≠ 0 → ∑ x, ((D x : ℤ) * (D (x + g) : ℤ)) = lam

def HasSDS (G : Type*) [AddCommGroup G] [Fintype G] (k : ℕ) (lam : ℤ) : Prop :=
  ∃ D : G → SignType, IsSDS k lam D
\end{Verbatim}

\noindent A single case then reads, for example
(\texttt{SignedDifferenceSets.Answer.N3.not\_hasSDS}),

\begin{Verbatim}[fontsize=\small]
theorem not_hasSDS : ¬ HasSDS (ZMod 3 × ZMod 3) 8 1
\end{Verbatim}

\noindent Every case's statement comes from its table entry, with the group
$\mathbb{Z}_{m_1} \times \dots \times \mathbb{Z}_{m_r}$ of its invariant factors
$[m_1, \dots, m_r]$: written out, as above, or through one definition \texttt{Grp} that unfolds to that product (Lean proves it has
$m_1 \cdots m_r$ elements). So besides the definition of an SDS, a reader needs to check only
that each statement's group, $k$ and $\lambda$ match its table entry, which a script can
compare.

\noindent From the repository root,

\begin{Verbatim}[fontsize=\small]
lake build --wfail SignedDifferenceSets
\end{Verbatim}

\noindent checks every result. For any result, \texttt{\#print axioms} shows only the three
standard axioms \texttt{propext}, \texttt{Classical.choice} and \texttt{Quot.sound}. No result
uses \texttt{native\_decide}: every computation is checked by Lean's kernel. A build from
scratch takes hours, and single modules need up to 12.3~GB of RAM; the problem's README says
how to build it on smaller machines.

\section{What is still open}\label{sec:open}

\stillOpen{} cases remain open in the database after our results: \csname openBelow50\endcsname{} with $v < 50$,
\csname openBelow100\endcsname{} with $50 \le v < 100$, \csname openBelow200\endcsname{} with
$100 \le v < 200$ and \csname openBelow500\endcsname{} with $200 \le v < 500$. Roughly, they
fall into four kinds: cases where a multiplier applies but the orbit search is too large;
cases where a rule above fails by a small margin; $p$-groups with $n$ a power of $p$, where a
construction seems more likely; and a core where no multiplier applies and every congruence
modulus is below $c$. We expect most of the core to have no SDS, but none of the local
arguments above can show it. Work on these cases continues; results beyond this summary,
including further corrections to the database, will be reported separately.

\section*{Use of AI}

I did this work with AI tools (Anthropic's Claude models). They searched the literature, found
and proved the nonexistence rules, ran the searches, wrote the Lean proofs and drafted this
note; I directed the work and decided which results to keep. So that nothing rests on the AI
alone, every case we settle or correct is checked by Lean~4, with no unproved steps and no
axioms beyond the three standard ones (the cyclic family above and the remark
on~\cite[Thm.~4.1]{he-chen-ge} are on paper), and before I accepted a result a separate AI agent tried to break it:
its proof, its Lean statement, the cases it settles and its credit to earlier work; these
reports are public too, in the problem's folder \texttt{reviews/}. The Lean statements are the
authoritative form of every settled case; all code, proofs and data are public (\repo).

\appendix
\section{The cases settled with a set}\label{app:sets}

\begingroup\small
\begin{longtable}{@{}llll@{}}
\toprule
Case & $(v,k,\lambda)$ & Group & Solution \\ \midrule
\endfirsthead
\toprule
Case & $(v,k,\lambda)$ & Group & Solution \\ \midrule
\endhead
\csname @@input\endcsname build/sets % TeX's primitive \input: LaTeX's breaks inside a table
\bottomrule
\end{longtable}
\endgroup

\bibliographystyle{amsplain}
\bibliography{refs}

\end{document}
